%% =============================================================================
%% VALIDATION
%% =============================================================================

\section{Numerical validation} \label{sec:validation}

Numerical accuracy is essential for scientific software. This section documents
the systematic validation of \pkg{prompt2analytics} against established
\proglang{R} packages, demonstrating that results match reference implementations
to high precision across all implemented methods.

%% -----------------------------------------------------------------------------
\subsection{Validation methodology} \label{sec:validation-methodology}

Our validation framework follows a rigorous methodology:

\begin{enumerate}
  \item \textbf{Reference selection}: For each method, we identify the canonical
    \proglang{R} implementation used in the econometrics and statistics
    literature (Table~\ref{tab:reference-packages}).

  \item \textbf{Test case design}: Each method is validated against multiple
    test cases:
    \begin{itemize}
      \item Synthetic data with known data-generating process (DGP)
      \item Real datasets with published results (e.g., Grunfeld, Longley)
      \item Edge cases (small samples, collinearity, unbalanced panels)
    \end{itemize}

  \item \textbf{Precision thresholds}: We define tolerance levels that vary
    by sample size and method type (Table~\ref{tab:tolerance-thresholds}).

  \item \textbf{Automated testing}: Validation tests are integrated into the
    continuous integration pipeline, ensuring results remain accurate as the
    codebase evolves.
\end{enumerate}

All validation tests use fixed random seeds (seed = 42) for reproducibility.
Reference \proglang{R} scripts are provided in the \code{validation/scripts/}
directory of the source repository.


%% -----------------------------------------------------------------------------
\subsection{Reference implementations} \label{sec:validation-reference}

Table~\ref{tab:reference-packages} lists the primary \proglang{R} packages
used for validation. We selected packages that are widely used in applied
econometrics and have established reputations for numerical accuracy.

\begin{table}[htbp]
\centering
\caption{Reference \proglang{R} packages for numerical validation}
\label{tab:reference-packages}
\begin{tabular}{llll}
\toprule
Category & Package & Version & Functions \\
\midrule
Regression & stats & 4.3.x & \fct{lm} \\
Robust SEs & sandwich & 3.0.x & \fct{vcovHC}, \fct{vcovCL} \\
Diagnostics & lmtest & 0.9.x & \fct{bptest}, \fct{dwtest} \\
Panel data & plm & 2.6.x & \fct{plm}, \fct{phtest} \\
HDFE & lfe & 2.8.x & \fct{felm} \\
IV/2SLS & AER & 1.2.x & \fct{ivreg} \\
Discrete choice & stats & 4.3.x & \fct{glm} (binomial) \\
Time series & vars & 1.5.x & \fct{VAR}, \fct{irf} \\
Forecasting & forecast & 8.21.x & \fct{Arima}, \fct{mstl} \\
Clustering & stats & 4.3.x & \fct{kmeans}, \fct{hclust} \\
\bottomrule
\end{tabular}
\end{table}

Secondary validation against \proglang{Python} packages (statsmodels,
linearmodels, scikit-learn) is performed for methods where multiple
reference implementations exist, providing additional confidence in
correctness.


%% -----------------------------------------------------------------------------
\subsection{Tolerance thresholds} \label{sec:validation-tolerance}

Numerical precision requirements vary by sample size and estimator type.
Table~\ref{tab:tolerance-thresholds} presents our tolerance guidelines.

\begin{table}[htbp]
\centering
\caption{Numerical tolerance thresholds by sample size}
\label{tab:tolerance-thresholds}
\begin{tabular}{lccc}
\toprule
Sample size & Coefficients & Standard errors & P-values \\
\midrule
$n < 100$ & $10^{-6}$ & $10^{-4}$ & 0.01 \\
$n = 100$--$1{,}000$ & $10^{-8}$ & $10^{-6}$ & 0.001 \\
$n > 1{,}000$ & $10^{-10}$ & $10^{-8}$ & 0.0001 \\
\midrule
Iterative methods (HDFE, MLE) & \multicolumn{3}{c}{$10^{-5}$ (convergence-dependent)} \\
\bottomrule
\end{tabular}
\end{table}

For iterative estimation procedures such as high-dimensional fixed effects
and maximum likelihood, slightly larger tolerances are applied due to
differences in convergence criteria between implementations.


%% -----------------------------------------------------------------------------
\subsection{Validation results} \label{sec:validation-results}

We present detailed validation results for core methods. All 70+ implemented
methods have corresponding validation tests in the test suite.

\subsubsection{Ordinary Least Squares}

OLS estimation is validated against the Longley dataset~\citep{Longley:1967},
a classic benchmark for numerical accuracy in regression software due to its
high multicollinearity.

\begin{table}[htbp]
\centering
\caption{OLS validation: Longley dataset ($n = 16$, $k = 6$)}
\label{tab:validation-ols}
\begin{tabular}{lrrr}
\toprule
Coefficient & \proglang{R} \fct{lm} & \pkg{prompt2analytics} & Difference \\
\midrule
(Intercept) & $-3482.259$ & $-3482.259$ & $< 10^{-8}$ \\
GNP.deflator & $0.01506$ & $0.01506$ & $< 10^{-8}$ \\
GNP & $-0.03582$ & $-0.03582$ & $< 10^{-8}$ \\
Unemployed & $-0.02020$ & $-0.02020$ & $< 10^{-8}$ \\
Armed.Forces & $-0.01033$ & $-0.01033$ & $< 10^{-8}$ \\
Population & $-0.05110$ & $-0.05110$ & $< 10^{-8}$ \\
Year & $1.82915$ & $1.82915$ & $< 10^{-8}$ \\
\midrule
$R^2$ & $0.9955$ & $0.9955$ & $< 10^{-4}$ \\
Adj.\ $R^2$ & $0.9925$ & $0.9925$ & $< 10^{-4}$ \\
\bottomrule
\end{tabular}
\end{table}

\subsubsection{Robust standard errors}

Heteroskedasticity-consistent standard errors (HC0--HC3) are validated against
the \pkg{sandwich} package~\citep{Zeileis:2004}. Table~\ref{tab:validation-hc}
shows achieved precision levels.

\begin{table}[htbp]
\centering
\caption{Robust standard error precision vs.\ \pkg{sandwich}}
\label{tab:validation-hc}
\begin{tabular}{lcc}
\toprule
HC type & Description & Precision vs.\ \proglang{R} \\
\midrule
HC0 & White (1980) & $< 10^{-10}$ \\
HC1 & Small-sample correction & $< 10^{-10}$ \\
HC2 & Leverage-adjusted & $< 10^{-8}$ \\
HC3 & MacKinnon-White (1985) & $< 10^{-8}$ \\
\bottomrule
\end{tabular}
\end{table}

The slightly reduced precision for HC2 and HC3 reflects numerical differences
in leverage (hat matrix diagonal) computation, which propagate through the
variance formula.

\subsubsection{Panel data estimation}

Fixed and random effects estimators are validated against the \pkg{plm}
package~\citep{Croissant+Millo:2008} using the Grunfeld investment dataset
($n = 200$, 10 firms $\times$ 20 years).

\begin{table}[htbp]
\centering
\caption{Panel FE validation: Grunfeld dataset}
\label{tab:validation-fe}
\begin{tabular}{lrrr}
\toprule
Statistic & \proglang{R} \pkg{plm} & \pkg{prompt2analytics} & Difference \\
\midrule
$\beta_{\text{value}}$ & $0.110124$ & $0.110124$ & $< 10^{-6}$ \\
$\beta_{\text{capital}}$ & $0.310065$ & $0.310065$ & $< 10^{-6}$ \\
SE($\beta_{\text{value}}$) & $0.011857$ & $0.011857$ & $< 10^{-5}$ \\
SE($\beta_{\text{capital}}$) & $0.017355$ & $0.017355$ & $< 10^{-5}$ \\
Within $R^2$ & $0.7668$ & $0.7668$ & $< 10^{-4}$ \\
F-statistic & $309.01$ & $309.01$ & $< 0.01$ \\
\bottomrule
\end{tabular}
\end{table}

The Hausman test for choosing between fixed and random effects produces
identical chi-squared statistics to \proglang{R}'s \fct{phtest} function
(see Section~\ref{sec:examples} for output).

\subsubsection{High-dimensional fixed effects}

Two-way fixed effects using the alternating projections (MAP) algorithm
\citep{Gaure:2013} are validated against the \pkg{lfe} package.

\begin{table}[htbp]
\centering
\caption{HDFE validation: Grunfeld dataset (firm + year FE)}
\label{tab:validation-hdfe}
\begin{tabular}{lrrr}
\toprule
Statistic & \proglang{R} \pkg{lfe} & \pkg{prompt2analytics} & Difference \\
\midrule
$\beta_{\text{value}}$ & $0.117716$ & $0.117716$ & $< 10^{-5}$ \\
$\beta_{\text{capital}}$ & $0.357916$ & $0.357916$ & $< 10^{-5}$ \\
Within $R^2$ & $0.7201$ & $0.7201$ & $< 10^{-4}$ \\
F-statistic & $217.44$ & $217.44$ & $< 0.01$ \\
\bottomrule
\end{tabular}
\end{table}

\subsubsection{Regression diagnostics}

Diagnostic tests (Jarque-Bera, Breusch-Pagan, Durbin-Watson, VIF) are
validated against \pkg{lmtest}, with all statistics matching to at least
$10^{-4}$. Per-test tolerances are documented in
Appendix~\ref{app:implementation}.


%% -----------------------------------------------------------------------------
\subsection{Test coverage summary} \label{sec:validation-coverage}

Table~\ref{tab:validation-summary} summarizes validation status across
method categories.

\begin{table}[htbp]
\centering
\caption{Validation test coverage by category}
\label{tab:validation-summary}
\begin{tabular}{lccl}
\toprule
Category & Methods & Tests & Reference \\
\midrule
Regression & 4 & 12 & stats, sandwich \\
Panel data & 4 & 8 & plm, lfe \\
Causal inference & 3 & 6 & AER \\
Discrete choice & 2 & 4 & stats \\
Time series & 5 & 10 & vars, forecast, urca \\
Machine learning & 7 & 14 & stats, dbscan \\
Diagnostics & 6 & 12 & lmtest \\
Statistical tests & 20+ & 40+ & stats \\
\midrule
\textbf{Total} & \textbf{51+} & \textbf{106+} & \\
\bottomrule
\end{tabular}
\end{table}

All validation tests are integrated into the continuous integration pipeline
and run automatically on each code change. The test suite can be executed
locally via:

\begin{Sinput}
cargo test -p p2a-core -- test_validate
\end{Sinput}


%% -----------------------------------------------------------------------------
\subsection{Known differences} \label{sec:validation-differences}

A small number of intentional differences exist between \pkg{prompt2analytics}
and reference implementations, including default robust SE type (HC1 vs.\
HC3), coefficient naming conventions, $R^2$ variant selection, and
convergence tolerances. These do not affect the validity of statistical
inference; details are in Appendix~\ref{app:known-differences}.

